% Source: sources/returns/S6_04_MAPS_MANIFOLDS/sections/manifolds.tex:2601–2842
\paragraph{Theorem D.2.13 (Proved: all smooth-fibre logs in four configurations).}
Consider the following four marked rational elliptic surfaces,
with $G$ a Mordell--Weil generator:
\[
\begin{array}{c|l|r}
\mathrm{OS}&\text{ordered fibres}&p\\\hline
45&I_8,I_1,I_1,I_1,I_1^\bullet&8\\
47&I_7,I_2,I_1,I_1,I_1^\bullet&14\\
55&I_5,I_4,I_1,I_1,I_1^\bullet&20\\
56&I_5,I_3,I_2,I_1,I_1^\bullet&30.
\end{array}
\]
Use $P=pG$, $Q=G$, and the degree-one linearization with zero
at $I_1^\bullet$. Keep the original four primitive one-node
fillings and the one-component hexagonal cusp. At an additional
smooth fibre, perform the logarithmic transformation of order
$m>0$ and integral translation vector
\[
 v=a e_1+b e_2+c\delta+n\gamma,\qquad
 \gcd(m,a,b,c,n)=1.
 \tag{D.S1}
\]
This is the exact freeness condition; no coprimality condition
on $m,n$ alone is imposed. Put
\[
 t=\gcd(m,n,a,b)\,\gcd(m,n,p,c).
 \tag{D.S2}
\]
The two factors in (D.S2) are coprime, and $t\mid m$.
For $n\ne0$, also $t\mid n$, and the result has
\[
\begin{aligned}
 \pi_1(X)&=H_1(X)=\mathbb Z/|n|,\qquad H_5(X)=0,\\
 H_2(X)&=H_3(X)=G_{m,n,t},\qquad H_4(X)=\mathbb Z/|n|,\\
 G_{m,n,t}&=\mathbb Z/t\oplus\mathbb Z/m
               \oplus\mathbb Z/(m|n|/t).
\end{aligned}
\tag{D.S3}
\]
For $n=0$, the answer is instead
\[
\begin{aligned}
 \pi_1(X)&=H_1(X)=H_5(X)=\mathbb Z,\\
 H_2(X)&=H_3(X)=\mathbb Z^2\oplus\mathbb Z/t\oplus\mathbb Z/m,
 \qquad H_4(X)=\mathbb Z^2.
\end{aligned}
\tag{D.S4}
\]
In both cases $H_0=H_6=\mathbb Z$ and $\chi(X)=2$.
The simply connected members are precisely $n=\pm1$; they
have $H_2=H_3=(\mathbb Z/m)^2$. Consequently the sphere
members in these four families are precisely $m=1$, $n=\pm1$.
For $m>1$ the theorem gives simply connected complex rational
homology spheres with two cyclic factors of order $m$ in their
middle torsion. Their oriented smooth type is identified in
D.2.14 below.
\begin{proof}
\emph{Geometric construction and complete affine scope.}
In these four configurations the height of $G$ is $1/p$.
The section $P$ is narrow at every fibre, whereas $Q$ has
unit component residue at each higher multiplicative fibre:
respectively $3$ modulo $8$; $(5,1)$ modulo $(7,2)$;
$(2,1)$ modulo $(5,4)$; and $(1,1,1)$ modulo $(5,3,2)$.
Thus $M=\mathcal O(P-O)$ has degree zero on each component,
and the quotient by $Q$ cycles through all components at each
weight-zero fibre. The geometric argument of D.2.12 identifies
each such filling with the primitive one-node model, including
all three normal tori. The mixed height is one, so
$M\otimes t_Q^*M^{-1}$ is the pullback of $\mathcal O(1)$.
Its chosen simple zero gives the stated hexagonal cusp.

Over a small disc at the additional regular value, the torus
family has a holomorphic origin and trivial integral local
system. Pull it back by $z=s^m$ and quotient by
\[
 (x,s)\longmapsto (x+v(s)/m,e^{2\pi i/m}s).
\]
Here $v(s)$ is the holomorphic period corresponding to the
integral vector $v$. The action is free exactly when (D.S1)
holds: away from $s=0$ its nonidentity powers move the base,
and on the central torus translation by $v/m$ has exact
order $m$ precisely under (D.S1). The usual logarithmic
collar, given by the translation
$(\log s)v(s)/(2\pi i)$, glues this proper quotient to the
original family; compare \cite[Proposition 2.3]{manViaclovskyLogs26}.
Changing $v$ by a multiple of $m$ retains its specified
integral collar clutch, rather than discarding that clutch.
The reduced central fibre is the torus with lattice
\[
 \Lambda'=\Lambda+\mathbb Z(v/m).
 \tag{D.S5}
\]
Its normal complex line is flat of finite order. Since the
integral $H^2$ of a torus is torsion-free, this line is
topologically trivial. Thus the boundary-to-fibre map is
onto on $H_2$, and its kernel consists of its four normal
circle tori.

\emph{Primitive global graph calculation.}
The marking is $(e_1,e_2,\delta,\gamma)$ and the invariant
coordinate is $\chi=\gamma^*$. Insert $A_0=I$ for the extra
regular fibre before the five original monodromies
$A_1,\ldots,A_5$. Their product is $I$.
All original vanishing lattices together generate precisely
the primitive $\ker\chi$; their combined matrix has three
unit Smith entries. Therefore van Kampen reduces the group
to the generator $\gamma$. All singular meridians vanish,
and the base relation kills the remaining regular meridian
$\xi$. The new relation $m\xi-v=0$ is then $n\gamma=0$.
This proves the assertion about the full fundamental group.

For completeness the integral degree-two calculation is
recorded explicitly. Set
\[
 D=[A_0-I\ A_1-I\ \cdots\ A_4-I]:\mathbb Z^{20}\to\mathbb Z^4.
\]
It has rank three and primitive image. As in D.2.12, denote
the five direct boundary embeddings into the graph chains by
$j_i$, and set
\[
 j_5(x)=-\bigl(A_1A_2A_3A_4x,\ A_2A_3A_4x,
                      \ A_3A_4x,\ A_4x,\ x\bigr).
\]
The unchanged fillings supply fourteen normal tori: three
at each of the four one-node fibres and two at the cusp.
If $B$ is their $20\times14$ graph matrix, then
\[
 K=[B\ j_0(e_1)\ j_0(e_2)\ j_0(\gamma)]
 \tag{D.S6}
\]
is a primitive integral basis of $\ker D$.
The retained integer matrix $L$ satisfies $LK=I_{17}$,
and the last three coordinates of $Lj_0(x)$ are
$(x_1,x_2,x_4)$. In particular the residual $\delta$ sweep
vanishes. These identities are verified for all four
constant marked monodromy tuples in the accompanying exact
program, which saves the matrices and their integer inverses.

The exterior coinvariants are free of rank one, detected by
\[
 \omega=p e_1^*e_2^*-\delta^*\gamma^*.
\]
Their presentation has five unit Smith entries. Every column
of (D.S6) is an actual fixed-loop sweep, on which the constant
vertical form $\omega$ has value zero. This specifies an
integral splitting of regular-family $H_2$.
After eliminating the fourteen unchanged normal tori,
the remaining group is freely generated by a fibre class
of $\omega$-value one and the three sweeps in (D.S6).

\emph{The exact normal map.}
The lattice (D.S5) is generated by $e_1,e_2,\delta,\gamma,g=v/m$,
with its single primitive relation
$mg=a e_1+b e_2+c\delta+n\gamma$.
For an integral vector $x$, its normal torus has fibre
coordinate $-\omega(x,v)$ and sweep $j_0(mx)$.
The $g$ normal torus has fibre coordinate zero and sweep
$j_0(v)$. One can read these formulas directly on the
affine boundary torus: its normal circle is $m\xi-v$,
so the first contribution is $-x\wedge v$, and the second
is $mx$ swept around $\xi$. This gives the complete
$4\times5$ matrix
\[
 C=\begin{pmatrix}
 -pb&pa&n&-c&0\\
 m&0&0&0&a\\
 0&m&0&0&b\\
 0&0&0&m&n
 \end{pmatrix}.
 \tag{D.S7}
\]
Its columns satisfy exactly the stated lattice relation,
so the five-generator presentation introduces no additional
normal torus.

\emph{Smith reduction for every integer choice.}
The greatest common divisor of all entries of $C$ is one,
by (D.S1). Its maximal minors, up to signs, are
\[
 m^3n,\quad m^2n^2,\quad m^2nc,\quad m^2nb,\quad m^2na.
\]
For $n\ne0$ their gcd is $m^2|n|$.
The lower determinantal divisors can be checked over each
local ring $\mathbb Z_{(\ell)}$, as follows.
If $m$ is a unit, three diagonal pivots leave just the
relation $n$.
If $n$ is a unit and $m$ is not, its pivot in the top row
and its pivot in the last column leave two copies of $m$.
In the remaining case $\ell\mid m,n$, condition (D.S1)
forces $c$ or one of $a,b$ to be a unit.
If $c$ is a unit, pivot on $-c$. Row operations give one
copy of $m$ together with the two-row presentation
\[
 \begin{pmatrix}m&0&h\\0&mn&n\end{pmatrix},
 \qquad (h)=(a,b)\text{ in }\mathbb Z_{(\ell)}.
\]
For $n\ne0$ the two nonzero determinantal divisors of this block
are $\gcd(m,n,a,b)$ and $m|n|$, up to local units.
For $n=0$ it has rank one and entry gcd $\gcd(m,a,b)$.
If instead $a$ or $b$ is a unit, pivot on that entry in
the last column. The remaining three-row presentation is
equivalent, up to units, to
\[
 \begin{pmatrix}p&n&-c\\m&0&0\\0&0&m\end{pmatrix}.
\]
Its first two determinantal divisors are
$\gcd(m,n,p,c)$ and $m\gcd(m,n,p,c)$.
The two gcd factors in (D.S2) have disjoint prime divisors:
a common divisor would divide $m,a,b,c,n$, contrary to
(D.S1). These local calculations therefore give the global
Smith entries
\[
 (1,t,m,m|n|/t)\quad(n\ne0),\qquad
 (1,t,m,0)\quad(n=0).
 \tag{D.S8}
\]
The same unit-pivot argument proves the $n=0$ formula
directly, without treating a finite group order as a limit.

The degree-one Mayer--Vietoris map has free source and
target of rank twenty-four. Its cokernel is the already
computed $H_1(X)$. Thus its kernel is zero when $n\ne0$,
and is a free copy of $\mathbb Z$ when $n=0$.
All local degree-two boundary maps are onto, including
the smooth logarithmic piece by (D.S5).
Consequently the full remaining degree-two cokernel is
$\operatorname{coker} C$, with precisely this extra free
relative term in the $n=0$ case. A free quotient splits
the resulting abelian-group extension. This proves $H_2$.
Only the hexagonal cusp contributes Euler characteristic,
giving two. Integral duality and the universal coefficient
theorem now give (D.S3)--(D.S4); in particular $b_3=2$
when $n=0$, and torsion appears in both $H_2$ and $H_3$.
The sphere conclusion follows from $\Theta_6=0$.

The program \texttt{code/manifolds\_semistable.py} verifies
all four constant primitive graph and exterior calculations
and 2924 full $18\times19$ Smith presentations, including
388 cases with $n=0$. The proof of (D.S8) has no bounds
on any parameter. The source search is recorded separately;
no novelty or priority claim is made.
\end{proof}
